\section{Prompts Used in \name}
\label{plan}
\subsection{Meta-agent}
\label{app_meta-agent}
The one-shot prompt for meta-agent's fast decomposition and allocation is shown in Figure \ref{meta_agent}.
\begin{figure}[H]
\centering
\begin{lstlisting}[frame=single, basicstyle=\scriptsize\ttfamily, breaklines=true]
meta_agent_prompt = '''
You are a planning agent. You are responsible for decomposing the given query into sub-tasks and choosing the most suitable agent for each sub-task. Your main goal is to efficiently and accurately complete task planning based on the descriptions of agents provided, ensuring the coherence and quality of the sub-tasks. 
Please output the sub-tasks and corresponding agents in the following JSON format: [{"task": task_description, "id": task_id, "name": name_of_agent, "reason": your_detailed_reason_for_the_choice, "dep":dependency_task_ids}]. In this format, "task" is the description of the sub-task, which will be used as the input of the chosen agent; "dep" denotes the id of the previous sub-task which generates a new resource relied by the current sub-task. 
The available agents and the corresponding descriptions are: [code_agent: Generate code in Python for precise computations to solve the given task. math_agent: Answer math questions by reasoning step-by-step. search_agent: Call Bing Search API to obtain information regarding the given task. commonsense_agent: Answer the given question using commonsense reasoning.]. 
---
Here is an example:
User query: If a plane can carry 300 passengers and decides to fly from China to Indonesia, how many full flights are needed to transport 1\% of the population of China to Indonesia?
Output:
[
    {
        "task": "Determine the population of China.",
        "id": 1,
        "name": "search_agent",
        "reason": "The search_agent can find the most recent and accurate population data for China.",
        "dep": []
    },
    {
        "task": "Calculate 1% of the population of China.",
        "id": 2,
        "name": "math_agent",
        "reason": "The math_agent can reason through the calculation step-by-step.",
        "dep": [1]
    },  
    {
        "task": "Determine the number of full flights needed to transport 1% of the population of China to Indonesia, given that each plane can carry 300 passengers.",
        "id": 3,
        "name": "math_agent",
        "reason": "The math_agent can reason through the division and rounding process.",
        "dep": [2]
    }
]
---
Given the user query, output the task plan with the format above directly. Make sure all the important information such as nouns or numbers are included in the sub-tasks. If you think the query can be done with just one agent, you can output only one sub-task.
'''
\end{lstlisting}
\caption{One-shot prompt for the meta-agent.}
\label{meta_agent}
\end{figure}

\subsection{Replan}
\label{app_replan}
The replan prompt is shown in Figure \ref{replan}.
\begin{figure}[H]
\centering
\begin{lstlisting}[frame=single, basicstyle=\scriptsize\ttfamily, breaklines=true]
replan_prompt = '''
You are a planning agent. You are preliminarily responsible for decomposing the given query into sub-tasks and choose the most suitable agent for each sub-task according to the following json format: [{"task": task_description, "id": task_id, "name": name_of_agent, "reason": your_detailed_reason_for_the_choice, "dep":dependency_task_ids}]. In this format, "task" is the description of the sub-task, which will be used as the input of the chosen agent; "dep" denotes the id of the previous sub-task which generates a new resource relied by the current sub-task. 
The available agents and the corresponding descriptions are: [code_agent: Generate code in Python for precise computations to solve the given task. math_agent: Answer math questions by reasoning step-by-step. search_agent: Call Bing Search API for obtaining information regarding the given task. commonsense_agent: Answer the given question using commonsense reasoning.]. 
Given the user query: %s, the preliminary task decomposition is: %s.
But the sub-task: %s cannot be solved by any agent. Now you are responsible for replaning this sub-task based on agents' capabilities. Output the new sub-task with the format above directly.
'''
\end{lstlisting}
\caption{Replan prompt.}
\label{replan}
\end{figure}

\subsection{Plan in Detail}
\label{app_plan_in_detail}
The prompt for planning in detail is shown in Figure \ref{planindetail}.
\begin{figure}[H]
\centering
\begin{lstlisting}[frame=single, basicstyle=\scriptsize\ttfamily, breaklines=true]
plan_in_detail_prompt = '''
You are a planning agent. You are preliminarily responsible for decomposing the given query into sub-tasks and choose the most suitable agent for each sub-task according to the following json format: [{"task": task_description, "id": task_id, "name": name_of_agent, "reason": your_detailed_reason_for_the_choice, "dep":dependency_task_ids}]. In this format, "task" is the description of the sub-task, which will be used as the input of the chosen agent; "dep" denotes the id of the previous sub-task which generates a new resource relied by the current sub-task. 
The available agents and the corresponding descriptions are: [code_agent: Generate code in Python for precise computations to solve the given task. math_agent: Answer math questions by reasoning step-by-step. search_agent: Call Bing Search API for obtaining information regarding the given task. commonsense_agent: Answer the given question using commonsense reasoning.]. 
Given the user query: %s, the preliminary task decomposition is: %s.
But the sub-task: %s cannot be solved only with agent %s. Now you are responsible for planning this sub-task in detail and choose the most suitable agents based on agents' capabilities. Output the new sub-tasks with the format above directly. Make sure that there are no duplicate content between new sub-tasks and the given preliminary task decomposition.
'''
\end{lstlisting}
\caption{Prompt for planning in detail.}
\label{planindetail}
\end{figure}

\subsection{Re-describe}
\label{app_redescribe}
The prompt for re-describing the sub-task is shown in Figure \ref{redescribe}.
\begin{figure}[H]
\centering
\begin{lstlisting}[frame=single, basicstyle=\scriptsize\ttfamily, breaklines=true]
redescribe_subtask_prompt = '''
Rewrite the following sentence based on the given example, while keeping the key information unchanged. Besides, output the rewritten sentence in the form like ***rewritten***.
---
Here is an example:
Example sentence: 'Determine the population of the United States in 2022.'
Sentence to be rewritten: 'Assess the population of China in 2022.'
Rewritten sentence: 'Determine the population of China in 2022.'
Output: ***'Determine the population of China in 2022.'***
---
Example sentence: %s
Sentence to be rewritten: %s
'''
\end{lstlisting}
\caption{Prompt for re-describing the sub-task.}
\label{redescribe}
\end{figure}

\subsection{Detector}
\label{app_detector}
The prompt for the detector is shown in Figure \ref{detector_prompt}.
\begin{figure}[H]
\centering
\begin{lstlisting}[frame=single, basicstyle=\scriptsize\ttfamily, breaklines=true]
plan_detector_prompt = '''
You are a plan detector responsible for analyzing the completeness and redundancy of the plan. Given the query and the plan formulated to solve the query, which involves several sub-tasks, you should do the following things:
1. **Detect whether the plan satisfies the completeness.**: Evaluate whether the set of subtasks covers all key aspects of the original task including important numbers and nouns. Specifically, check if each important element and requirement from the original task is addressed by at least one subtask. Provide a brief explanation if any key information is missing.
2. **Detect whether the plan satisfies the non-redundancy.**: Evaluate whether any two sub-tasks contain identical information and requirements. If there is any redundant part, list and provide suggestions for optimizing the plan.
---
For example:
Task: If a plane can carry 300 passengers and flies from Brazil to Nigeria with a full load, then returns with only 75% capacity filled, how many passengers in total has it transported between the two countries in one round trip?
Subtask 1: Determine the number of passengers transported from Brazil to Nigeria in one flight with a full load.    Dependency: []
Subtask 2: Determine the number of passengers transported from Nigeria to Brazil in one flight with 75% capacity filled.    Dependency: []
Subtask 3: Calculate the total number of passengers transported between Brazil and Nigeria in one round trip.    Dependency: [1, 2]
Analyse: This plan does not satisfy completeness because the subtask loses the information of 'a plane can carry 300 passengers' of the original task. This plan satisfies non-redundancy because each subtask has a unique focus and there is no overlap in the information covered.
Suggestions: Add the information of 'a plane can carry 300 passengers' to subtask 1 and subtask 2.
---
If there is no need to modify the plan, just return 'The plan satisfies completeness and non-redundancy.'.
'''
\end{lstlisting}
\caption{Prompt for the detector.}
\label{detector_prompt}
\end{figure}

\subsection{Evaluation}
The prompt for comparing the response with the ground truth is shown in Figure \ref{evaluate}.
\label{evaluation}
\begin{figure}[H]
\centering
\begin{lstlisting}[frame=single, basicstyle=\scriptsize\ttfamily, breaklines=true]
evaluate_prompt = '''
You are CompareGPT, a machine to verify the correctness of predictions. Answer with only yes/no.
You are given a question, the corresponding ground-truth answer and a prediction from a model. Compare the "Ground-truth answer" and the "Prediction" to determine whether the prediction correctly answers the question. The prediction may contain extra information, but a correct prediction includes the ground-truth answer. You can answer "yes" if the prediction includes the ground-truth answer. You must answer "no" if there are any specific details in the ground-truth answer that are not mentioned in the prediction. If the prediction states something as a possibility, treat it as a definitive answer. Note that the error within three decimal places is negligible.
---
Question: %s
Ground-truth answer: %s
[Start of the prediction]
%s
[End of the prediction]
'''
\end{lstlisting}
\caption{Prompt for comparing the response with the ground truth.}
\label{evaluate}
\end{figure}

\subsection{Agents}
The prompts for the code, math, search and commonsense agent are shown in Figure \ref{code}, \ref{math}, \ref{search}, \ref{commonsense} separately. Specifically, the commonsense agent can generate commonsense information like the metal melting points and the atomic numbers of helium and hydrogen.
\label{agents}
\begin{figure}[H]
\centering
\begin{lstlisting}[frame=single, basicstyle=\scriptsize\ttfamily, breaklines=true]
code_agent_prompt = '''You ara a code agent. You can be used for : 1) computing large numbers, fractions or decimals. 2) counting or averaging long lists of numbers. 3) performing date-related operations, such as counting the number of days between two dates. Write code in Python to solve the given task with history. Give the code in the following form directly. 
- Here is an example: 
Task: Calculate the combined population of China and India in 2022.
History: The answer of 'Determine the population of China in 2022' is 1.412B. The answer of 'Determine the population of India in 2022' is 1.417B.
Code:
```python
# Given populations
population_china_2022 = 1.412 * 10**9  # 1.412 billion
population_india_2022 = 1.417 * 10**9  # 1.417 billion

# Calculate combined population
combined_population_2022 = population_china_2022 + population_india_2022

# Print the result
print(f"The combined population of China and India in 2022 is {combined_population_2022} people.")
```
---
Task: %s
History: %s
Code:
'''

rewrite_code_agent_prompt = '''You are a rewrite agent. Given the input question, the code addressing this question and the corresponding output, rewrite the output into a complete sentence that integrates information from the question and the code output.
---
Question: %s
Code: %s
Code output: %s
Output:
'''
\end{lstlisting}
\caption{Prompt for the code agent.}
\label{code}
\end{figure}

\begin{figure}[H]
\centering
\begin{lstlisting}[frame=single, basicstyle=\scriptsize\ttfamily, breaklines=true]
math_agent_prompt = '''You are a math agent. You can answer math questions by reasoning step-by-step with the data provided in the question and history. Present the answer "ANS" to the subquestion in LaTeX using the format 'The answer is \boxed{ANS}.' without any units in the box.
---
Question: %s
History: %s
Solution: 
'''
\end{lstlisting}
\caption{Prompt for the math agent.}
\label{math}
\end{figure}

\begin{figure}[H]
\centering
\begin{lstlisting}[frame=single, basicstyle=\scriptsize\ttfamily, breaklines=true]
search_agent_prompt = '''You are a search agent. Write a concise, informative Bing Search query for obtaining information regarding the given task.
- Here is an example:
Task: Determine the population of China in 2022.
History: None
Search query: China population 2022
---
Task: %s
History: %s
Search query: 
'''
rewrite_search_agent_prompt = '''You are a rewrite agent. Given the search question, the response in the web_pages from the bing search api, answer the search question with the information from the response in concise words. Remove redundant information that is irrelevant to the question.
---
Question: %s
Answer_box: %s
Answer:
'''
\end{lstlisting}
\caption{Prompt for the search agent.}
\label{search}
\end{figure}

\begin{figure}[H]
\centering
\begin{lstlisting}[frame=single, basicstyle=\scriptsize\ttfamily, breaklines=true]
commonsense_agent_prompt = '''You are a commonsense agent. You can answer the given question with logical reasoning, basic math and commonsense knowledge.
---
Question: %s
History: %s
Solution: 
'''
\end{lstlisting}
\caption{Prompt for the commonsense agent.}
\label{commonsense}
\end{figure}


\section{Fault Tolerance: Unexpected Changes in Agent Availability}
\label{app: fault_tolerance_agent_availability}

We categorize changes in agent availability into two situations: those occurring during non-execution periods and those during execution.
\begin{itemize}
    \item \textit{During non-execution periods}: \name allows for the addition or removal of agents. The meta-agent can first broadcast a simple sync signal to determine agent availability. Only available agents are provided to the meta-agent for task allocation. 
    \item \textit{During execution}: Changes in agent availability during execution indicate that an agent selected for a task may unexpectedly become unavailable, leading to the meta-agent not receiving a response to this sub-task. In such scenarios, the meta-agent is required to reassign the sub-task to another agent with similar capabilities or to further decompose the task (i.e., plan-in-detail).
\end{itemize}
\section{Constructing Training Dataset}
The prompt for the scorer is shown in Figure \ref{scorer_prompt}. The criteria for manual scoring are consistent with the standards outlined in the prompt provided to the scorer.
\label{training}
\begin{figure}[H]
\centering
\begin{lstlisting}[frame=single, basicstyle=\scriptsize\ttfamily, breaklines=true]
scorer_prompt = '''
Please act as an impartial judge and evaluate the quality of the response provided by the %s to the user task. Your evaluation should consider three factors: correctness, relevance and completeness. Assign a score of 0, 1 or 2 for each factor and provide a brief explanation for your score. The following is the grading criteria.
---
Correctness
0: The response contains severe errors and is completely inaccurate.
1: The response has some errors, but the main content is generally correct
2: The response is completely accurate and fully meets the requirements of the task.
Relevance
0: The response is minimally relevant to the task and completely off-topic.
1: The response is somewhat relevant to the task but may include some unrelated content.
2: The response is highly relevant to the task, directly addressing the core issue without any unrelated content or deviation.
Completeness
0: The response lacks necessary detail or key information, resulting in an incomplete understanding of the task.
1: While the response addresses part of the task, more information or content is needed for completeness.
2: The response provides comprehensive information and detailed explanations.
---
Besides, summarize the final result in the form like '**Correctness: score, Relevance: score, Completeness: score**' at the end of your response, where score can be chosen from 0, 1 and 2.
---
Task
%s
[The Start of Agent's Response]
%s
[The End of Agent's Response]
'''
\end{lstlisting}
\caption{Prompt for the scorer.}
\label{scorer_prompt}
\end{figure}



We design a mapping function that converts different combinations of correctness, relevance, and completeness values into a final score, which is shown in Figure \ref{mapping}. Then, a score of 5 or higher is considered sufficiently high, while a score of 1 or lower is considered sufficiently low.
\begin{figure}[H]
\centering
\begin{lstlisting}[frame=single, basicstyle=\scriptsize\ttfamily, breaklines=true]
score_map = {
    (2, 2, 2): 8,
    (2, 1, 2): 7,
    (2, 2, 1): 6,
    (2, 1, 1): 5,
    (1, 2, 2): 4,
    (1, 1, 2): 3,
    (1, 2, 1): 2,
    (1, 1, 1): 1,
}
def level_score(correctness, relevance, completeness):
    return score_map.get((correctness, relevance, completeness), 0)
\end{lstlisting}
\caption{Mapping function.}
\label{mapping}
\end{figure}

\section{Experiments on More Datasets}
\subsection{Performance on DROP and IIRC}
\label{app:exp_more_datasets}
In addition to Husky-QA~\citep{kim2024husky}, we also conduct experiments on DROP~\citep{DROP} and IIRC~\citep{IIRC}. The experimental results are shown in Table~\ref{table:more_dataset}, which indicate that \name significantly outperforms the baselines, achieving at least a 5.0\% improvement on both DROP and IIRC. These experimental results further confirm the effectiveness and advancements of \name.
\begin{table}[h]
\centering
\caption{Accuracy (\%) comparisons between \name and baselines on DROP and IIRC.}
\label{table:more_dataset}
\begin{tabular}{lcc}
\toprule
Method & DROP & IIRC \\ \midrule
GPT-4o & 23.0 & 33.0 \\
CoT & 26.0 & 35.0 \\
Zero-Shot CoT & 24.5 & 33.0 \\
Meta-Agent & 25.5 & 32.5 \\
Meta-Agent: Traversal & 27.5 & 34.0 \\
REACT & 29.0 & 36.0 \\
HUSKY & 28.0 & 36.5 \\
\midrule
\name & \textbf{34.0} & \textbf{41.5}  \\
\name (reward model trained on Husky-QA) & \underline{32.0} & \underline{39.0} \\
\bottomrule
\end{tabular}
\end{table}


\subsection{Generalization Capability of the Reward Model}
\label{app:exp_generalization_reward_model}
To evaluate the generalization capability of the reward model, we train a reward model based on Husky-QA and utilize it in experiments on the DROP and IIRC datasets. As shown at the bottom in Table~\ref{table:more_dataset}, the reward model trained on one dataset demonstrates good generalization to other datasets (though it experiences a slight performance drop compared to the specifically trained reward model), achieving superior performance compared to baselines.





\section{Further Investigation of Expert Agents}
\label{app:further_expert_agents}
\subsection{Expert Agents}
\label{app:exp_expert_agents}
To investigate the effectiveness of expert agents powered by LLMs fine-tuned on task-specific datasets, we adopted Qwen2-Math-7B~\citep{yang2024qwen2} as the backbone LLM for the math agent and DeepSeek-Coder-V2~\citep{zhu2024deepseek} as the backbone LLM for the code agent. The experimental results show that integrating these expert agents results in an additional 1.8\% improvement on Husky-QA compared with using GPT-4. These results demonstrate the potential for further enhancements of \name by incorporating more powerful expert agents.

\subsection{Multiple Agents with Same Expertise}
\label{app:exp_similar_agents}
Based on \name, we set up an experiment involving 4 different math agents, using GPT-3.5, GPT-4o, Qwen2-Math-7B, and Llama-3.2-3B, respectively. We instruct the multi-agent system to solve the complex math problem from MATH~\citep{MATH}. 
With \name, the queries are decomposed into multiple queries and these agents run in parallel to resolve them.
The experimental results are shown in Table~\ref{table:exp_similar_agents}, demonstrating the effectiveness (at least 6\% improvements) of \name.
\begin{table}[h]
\centering
\caption{Accuracy (\%) comparisons on MATH.}
\label{table:exp_similar_agents}
\begin{tabular}{lccccc}
\toprule
 & GPT-3.5 & GPT-4o & Qwen2-Math-7B-Instruct & Llama-3.2-3B & \name  \\ \midrule
MATH & 43 & 62 & 66 & 36 & \textbf{72}\\
\bottomrule
\end{tabular}
\end{table}


\section{Additional Experimental Results and Discussions}
\label{app:additional_exp}

\subsection{Quantitative Evaluation on Solvability, Completeness, and Non-redundancy}
\label{app:quantitative_eval}
We conduct a quantitative evaluation to demonstrate the effectiveness of \name in terms of solvability, completeness, and non-redundancy. Specifically, we compare the decomposed sub-tasks provided by \name to those provided by GPT-4o. These sub-tasks are assessed by an LLM-based evaluator, providing binary scores for solvability, completeness, and non-redundancy, with scores of 1 for those that meet these principles and 0 for not. The experimental results (averaged scores), as summarized in Table~\ref{table:quantitative_eval}, show that \name achieves significant improvements.

\begin{table}[h]
\centering
\caption{Quantitative Evaluation on Solvability, Completeness and Non-redundancy.}
\label{table:quantitative_eval}
\begin{tabular}{lccccc}
\toprule
 & Solvability & Completeness & Non-redundancy  \\ \midrule
GPT-4o & 0.763 & 0.822 & 0.986 \\ 
\name & 0.938 & 0.969 &  0.993 \\
\bottomrule
\end{tabular}
\end{table}

\subsection{Effect of Different Representation Approaches for Agent Description}

In this study, we employ a combination of predefined natural language descriptions and representative works as the representations of agents’ descriptions. A natural language description can be manually provided or automatically generated, detailing the general and task-independent abilities of agents. While providing a comprehensive and detailed natural language description can be beneficial, it also requires effective prompt engineering. 
The representative (please refer to Sec.~\ref{subsec:task_modifications} for more details) works complement the natural language descriptions and are often task-dependent, allowing for continuous updates during execution.

We conducted experiments on Husky-QA to compare the effectiveness of different representation approaches. The results, shown in Table~\ref{table:exp_agent_descrption}, indicate that using natural language descriptions achieves significantly superior performance compared to using only representative works, which motivates the majority of existing studies to adopt natural language descriptions. Besides, incorporating task-specific representative works on top of natural language descriptions leads to a further 11.9\% performance boost, demonstrating the effectiveness of our proposed combined representation approach.

\begin{table}[h]
\centering
\caption{Comparisons among Natural Language Descriptions and Representative Works.}
\label{table:exp_agent_descrption}
\begin{tabular}{lc}
\toprule
Approaches & Accuracy (\%) \\ \midrule
Natural Language Descriptions & 31.8 \\ 
Representative Works & 16.8 \\
Both & 43.7 \\
\bottomrule
\end{tabular}
\end{table}

Besides, in this study, we employ the natural language descriptions adapted from Husky~\cite{kim2024husky} to ensure fair comparisons. To further investigate the effect of natural language descriptions, we provide two different versions: (i) LLM-generated: We instruct the LLM to generate the natural language descriptions based on how we describe the agents in Section~\ref{subsec:exp_settings}; (ii) Expert-written: The natural language descriptions crafted by a human expert. The comparisons on Husky-QA are shown in Table~\ref{table:exp_agent_descrption_diff}. From the table, we can observe that using simple, fully automated natural language descriptions may somewhat affect the overall performance. Besides, incorporating more expert insights into the generation of natural language descriptions can provide an additional performance boost. These experimental results align with the current intuitive understanding of prompt engineering.
\begin{table}[h]
\centering
\caption{Comparisons among Different Natural Language Descriptions.}
\label{table:exp_agent_descrption_diff}
\begin{tabular}{lc}
\toprule
 & Accuracy (\%) \\ \midrule
Descriptions adapted from Husky & 43.7 \\ 
LLM-generated descriptions & 40.4 \\
Expert-written descriptions & 44.2 \\
\bottomrule
\end{tabular}
\end{table}





\subsection{The Effectiveness of Detector in Handling Missing Dependencies}
\label{app:exp_detector}
We conduct an experiment to evaluate the effectiveness of the detector in handling missing dependencies. To be more specific, the detector is instructed to determine whether there are any missing dependencies before a sub-task is assigned to an agent for execution. If such dependencies exist, the execution results of these dependent sub-tasks must also be provided as inputs. The adopted prompt is provided in Figure~\ref{det_dep}. If there exists any missing information, the meta-agent will add sub-tasks to complete the dependencies.
The experiments conducted on Husky-QA show that, by enabling the detector to handle missing dependencies, \name obtains a 2.5\% performance increase in accuracy (from 43.7\% to 46.2\%). 

\begin{figure}
\centering
\begin{lstlisting}[frame=single, basicstyle=\scriptsize\ttfamily, breaklines=true]
dep_detect_prompt = '''
You are an intelligent detector tasked with determining whether the provided dependency information is sufficient to complete a given task. If the dependency information is insufficient, you will review all historical data to find supplemental information. If neither the dependency information nor the historical data is sufficient, you will identify and list the missing information. Besides, the information in the task or a given query can be viewed as available in the dependency information.

Input:
Task: {Task description}
Dependency Information: {Dependency information provided}
Historical Data: {Historical data (listed with serial numbers)}
Available query: {The query whose information can be viewd as available inthe dependency information.}

Requirements:
1. Assess Dependency Information: Evaluate if the provided dependency information is sufficient to complete the task. If sufficient, answer ***Yes***; if not, answer ***No***. Besides, if the answer is ***Yes***, there is no need to check the following requirements.
2. Review Historical Data: If the answer above is ***No***, check the historical data to identify any relevant entries that can supplement the task requirements. List the serial numbers of the relevant entries as ~~~numbers~~~.
3. Identify Missing Information: If the historical data also cannot supplement the missing details, explicitly list what information is still required to complete the task as $$$required additional information$$$.

Output Format (strictly follow this structure):
1. Is the dependency information sufficient: ***Yes***/***No***
2. Relevant information from historical data: ~~~numbers~~~ or ~~~None~~~
3. Missing information: $$$Specific missing information$$$ or $$$None$$$

Examples:
---
Input:
Task: Calculate the total population of China and the United States in 2022.
Dependency Information: 
China's population in 2022 is 1.4 billion. 
Historical Data:
1. China's population in 2022 is 1.4 billion. 
2. The United States population in 2022 is 330 million.
3. Total world population in 2022: 8 billion.
Available query: If the populations of China and India were combined in 2022, how many countries with a population of 70,850,000 each could be formed from this total population without leaving anyone out?
Output:
1. Is the dependency information sufficient: ***No***
2. Relevant information from historical data: ~~~2~~~
3. Missing information: $$$None$$$
---
Input:
Task: Calculate the total population of China and the United States in 2022.
Dependency Information: 
China's population in 2022 is 1.4 billion. 
Historical Data:
1. China's population in 2022 is 1.4 billion. 
2. Total world population in 2022: 8 billion.
Available query: If the populations of China and India were combined in 2022, how many countries with a population of 70,850,000 each could be formed from this total population without leaving anyone out?
Output:
1. Is the dependency information sufficient: ***No***
2. Relevant information from historical data: ~~~None~~~
3. Missing information: $$$The United States population in 2022$$$
----
Input:
Task: %s
Dependency Information: 
%s
Historical Data:
%s
Available query: %s
Output:
'''
\end{lstlisting}
\caption{Prompt for the detector in handling missing dependencies.}
\label{det_dep}
\end{figure}









\subsection{The Limitations on the Number of Agents}
\label{app:number_of_agents}
\name has no inherent limitations on the number of agents. However, as the meta-agent is powered by LLMs, increasing the number of agents is subject to the contextual window length restrictions of LLMs (descriptions of a large number of agents could exceed the LLM's output length limit, such as 128K tokens), and the effectiveness of LLMs could be affected by the ability to handle long contexts.

In practical applications, it is important to consider that as the number of agents increases, there may be redundant and functionally similar agents, which motivates us to apply some extension strategies to handle. For example, agents can be grouped according to their abilities. When the meta-agent allocates sub-tasks, it can first select an agent group for each sub-task and then further choose the most suitable agent within the group. This strategy respects the LLM's context window limits and enhances the accuracy of agent selection.






